\section{What the Stack Actually Is}
\label{sec:ble-what}

Before writing a line of binding code, three questions: what is in the binary, how do you talk to it,
and what does it expect you to do.

\subsection{Three Libraries, One Protocol}

\begin{table}[htbp]
    \centering
    \footnotesize
    \begin{tabular}{@{}p{0.34\textwidth}p{0.26\textwidth}p{0.34\textwidth}@{}}
        \toprule
        \textbf{Library} & \textbf{What it is} & \textbf{Size on this part} \\
        \midrule
        \codeinline{libbluetooth\_controller} & link layer: radio, timing, packets & ~600 KB \\
        \codeinline{libbluetooth\_host} & GAP, GATT, SM: the protocol logic & ~700 KB \\
        \codeinline{libbgcommon} & shared buffers, BGAPI parsing & ~40 KB \\
        \bottomrule
    \end{tabular}
    \caption{The three blobs of a Gecko SDK 4.x Bluetooth build. There is no fourth: the profiles and
    services you need are your code.}
    \label{tab:ble-blobs}
\end{table}

Together they are larger than most applications on this part, which is why a project chooses
Bluetooth deliberately rather than by default. All three are MSLA licensed and therefore not in this
book's repository; a clean checkout compiles the bindings and fails at link with a list of missing
symbols, which is the correct behaviour.

\subsection{BGAPI: Commands In, Events Out}

The API is asynchronous and message-based, and both halves matter:

\begin{minted}{text}
command:  sl_bt_advertiser_set_data(handle, ...)   ->  sl_status_t (immediate)
reply:    sl_bt_evt_advertiser_set_data_complete   (later, in the event queue)
\end{minted}

Every command returns a status that says ``accepted'' or ``refused'' — never ``done''. The completion
arrives as an event, and an application that does not drain the event queue eventually stops being
able to send commands, because the queue is finite and the stack does not grow it.

\begin{minted}{rust}
loop {
    if let Some(event) = ble.poll_event() {
        match event { /* Connected, Disconnected, GattWrite, ... */ }
    }
}
\end{minted}

That loop is not optional and it is not decoration: it is the other half of every call.

\subsection{What the Stack Assumes You Have Done}

The blobs are written for a vendor project, and they assume the surrounding C world exists:

\begin{itemize}
    \item \textbf{Clocks.} The radio needs the HFXO; the stack enables it, but if the crystal never
          settles the boot event never arrives.
    \item \textbf{Critical sections.} The CMSIS \codeinline{CORE\_*} helpers, which in a vendor
          project come from headers as \codeinline{static inline}.
    \item \textbf{An event loop.} Yours.
    \item \textbf{Memory.} Fixed buffers inside the blobs, sized by the stack configuration in the
          header. A Rust project cannot change them without regenerating that configuration.
    \item \textbf{Crypto} — for any connectable, bondable device, a PSA implementation the blob calls
          but does not contain.
\end{itemize}

None of that is in the API documentation, and all of it appears as a link error or a timeout on first
use. Section~\ref{sec:ble-symbols} is about the link errors; section~\ref{sec:ble-beacon} is about the
timeouts.
